\section{讨论与未来方向 \label{sec:discussion}}

第~\ref{sec:results} 节的结果清楚表明：
可以用只激发每时间片上极少数扩展自由度的算符来研究轻强子。
产生一个强子所需的受限模式集意味着：构造关联函数所需的
夸克传播的全部元素都可以直接测得。这并不妨碍用随机方法
\cite{Foley:2005ac} 来降低计算成本。

在本文描述的直接方案中，蒸馏空间由时间片上格点三维拉普拉斯算符的
最低本征矢量构造。这一定义并非唯一，
对蒸馏空间的更多不同选择做试验可能有用。
另外，本工作中使用的算符在夸克自旋指标上是对角的。
蒸馏空间的其他选择可以包括非平庸的自旋结构，
例如使用三维狄拉克算符最低模式张成的向量空间。

回顾一下，在本文检验的最简单实现中，当 $N=M$ 时
蒸馏算符成为恒等算符，不做任何涂抹。
这显然是不理想的特性；对方法做显而易见的改造——
给各本征模赋予不同的权重——即可规避。这个更一般的蒸馏算符为
\begin{equation}
        {\cal J}(t) = \sum_{k=1}^N v^{(k)}(t) \, f\big(\lambda_k(t)\big)
              \, v^{(k)\dag}(t).
\end{equation}
当采用许多本征矢量时，在本征矢量上加入权重函数
可用于更快地压制剧烈起伏的贡献。
使用指数权重会给出对式~\ref{eqn:smearing} 传统涂抹算法
越来越准确的近似。

由于正如第~\ref{subsec:results-wavefn} 节数据所示，光滑的蒸馏模式
几乎与截断尺度动力学完全退耦，
因此预期在固定物理体积下趋向连续极限时不需要增加本征矢量数目。
这一预期尚待数值检验。
该方法一个重大的潜在缺点是：随三维格点体积增大，
需要提高蒸馏算符的秩。求解狄拉克方程成本的相应线性增长
导致谱学算法具有 ${\cal O}(V^2)$ 的标度。
从我们给出的结果来看，随体积增大而必需的秩的增加同时也带来
统计涨落的下降。体积标度问题的可能解决方案是把该方法与
适当设计的随机估计方案相结合。这正在研究之中。

注意该方法并不能直接给出夸克传播子的所有元素，
只能给出与低能谱学相关的那些元素。在涉及同位旋标量算符插入的计算中
这会成为限制，因为那里需要具有夸克场全部动量分量的不连通圈图。
这类计算的一个著名例子是核子奇异含量的计算。
